\chapter*{Resumen}
\addcontentsline{toc}{chapter}{Resumen/Abstract\vspace{-0.5cm}}

Las Amenazas Persistentes Avanzadas corresponden a campañas de ataque dirigidas y sigilosas que se desarrollan en múltiples etapas a lo largo de meses, lo que les permite evadir los mecanismos de detección tradicionales. Estos mecanismos evalúan eventos de forma aislada y no la secuencia temporal completa que caracteriza a este tipo de campañas, situación que se ve agravada por el severo desbalance entre tráfico legítimo y malicioso presente en los escenarios reales de operación.

Este trabajo propone la implementación y validación de un modelo basado en Cadenas Ocultas de Markov para la detección de este tipo de amenazas a partir de tráfico de red, representando las etapas del ataque como estados no observables que se infieren desde la evidencia disponible. La investigación se enmarca en la metodología de Design Science Research y se sustenta en una revisión sistemática de la literatura, que fundamenta tanto la selección del enfoque probabilístico como el diseño del modelo y de los conjuntos de datos empleados para su evaluación.

La revisión sistemática, ejecutada sobre Scopus, Web of Science e IEEE Xplore, conformó un corpus de 77 estudios primarios y permitió identificar cuatro brechas: la dependencia de alertas previamente correlacionadas, la escasa aplicación de estos modelos directamente sobre tráfico de red, la falta de un conjunto de datos de referencia consolidado y el predominio de métricas globales que ocultan el desempeño sobre la clase minoritaria. A partir de ellas se diseñó un modelo cuyos estados se anclan a las etapas etiquetadas de los conjuntos DAPT2020 y Unraveled, que trata el desbalance sin alterar la estructura temporal de las secuencias y que se evalúa mediante métricas por etapa y de detección temprana, en comparación con Random Forest y XGBoost bajo las mismas condiciones experimentales. Un prototipo implementado en Python y evaluado de forma preliminar sobre ambos conjuntos muestra que la estimación supervisada conserva la correspondencia entre estados y etapas que el entrenamiento no supervisado pierde, mientras que con la configuración actual los modelos de referencia superan al modelo propuesto en la clasificación por ventana y en la detección a igual tasa de falsas alarmas. Estos resultados orientan los ajustes del diseño, como el análisis de la campaña a nivel de red y la selección del tamaño de ventana y de la discretización.

{\bf Palabras clave:} amenazas persistentes avanzadas, cadenas ocultas de Markov, detección temprana, detección de intrusiones, revisión sistemática de literatura, desbalance de clases


\clearpage
\chapter*{Abstract}
Advanced Persistent Threats are targeted, stealthy attack campaigns that unfold across multiple stages over several months, allowing them to evade traditional detection mechanisms. These mechanisms assess events in isolation rather than the full temporal sequence that characterises such campaigns, a limitation compounded by the severe imbalance between legitimate and malicious traffic found in real operating environments.

This work proposes the implementation and validation of a Hidden Markov Model for detecting this type of threat from network traffic, representing attack stages as unobservable states inferred from the available evidence. The research follows the Design Science Research methodology and is grounded in a systematic literature review, which supports the selection of the probabilistic approach as well as the design of the model and the datasets used for its evaluation.

The systematic review, conducted on Scopus, Web of Science and IEEE Xplore, yielded a corpus of 77 primary studies and revealed four gaps: reliance on previously correlated alerts, the limited application of these models directly to network traffic, the lack of a consolidated benchmark dataset, and the prevalence of global metrics that conceal performance on the minority class. Building on these gaps, the proposed model anchors its states to the labelled stages of the DAPT2020 and Unraveled datasets, handles class imbalance without altering the temporal structure of the sequences, and is evaluated with per-stage and early-detection metrics against Random Forest and XGBoost under the same experimental conditions. A Python prototype, evaluated in a preliminary way on both datasets, shows that supervised estimation preserves the correspondence between states and stages that unsupervised training loses, whereas with the current configuration the baseline models outperform the proposed model in window-level classification and in detection at equal false alarm rates. These results guide the design adjustments, such as a network-level analysis of the campaign and the selection of the window size and the discretization.

{\bf Keywords:} advanced persistent threats, hidden Markov models, early detection, intrusion detection, systematic literature review, class imbalance
